\begin{resumo}
\begin{SingleSpace}
O uso de \textit{smartphones} tem se intensificado ao longo dos anos, tornando-se parte essencial da vida cotidiana. Apesar das vantagens da conectividade, o uso excessivo desses dispositivos está associado a impactos físicos, mentais, cognitivos e sociais. Entre os principais problemas relatados na literatura estão distúrbios do sono, ansiedade, dificuldades de atenção, dores no pescoço e dores de cabeça. As soluções existentes, como bloqueios nativos de aplicativos, são fáceis de burlar pelos próprios usuários. Este trabalho propõe o desenvolvimento do aplicativo Desconecta, com abordagem gamificada, no qual grupos de usuários participam de desafios semanais para manter menor tempo de uso de tela em um determinado período. O aplicativo permite acompanhar o tempo de uso, definir metas e comparar resultados entre os participantes. O desenvolvimento incluiu pesquisa com usuários para coleta de dados e levantamento de requisitos, seguida da implementação da solução e de sua avaliação com o público-alvo. Nos testes realizados, observou-se redução no tempo médio de uso entre os participantes e bom engajamento na dinâmica proposta, indicando que a abordagem coletiva pode contribuir para um uso mais consciente e saudável da tecnologia.
\end{SingleSpace}
\vspace{\onelineskip}
\textbf{Palavras-chave}: uso excessivo de \textit{smartphones}; tempo de tela; saúde mental; \textit{design} persuasivo; aplicativo de controle

\end{resumo}

% Palavras-chave separadas e finalizadas por ponto